%% Heyting algebras, IPC, and prime-filter Kripke completeness.

\subsection{Heyting algebras and intuitionistic completeness}
\label{sec:p10-heyting-completeness}

\begin{hbdefinition}{def:p10-heyting-algebra}{Heyting algebra}
A Heyting algebra is a nonempty bounded distributive lattice
\(H=(H,\wedge,\vee,0,1)\), equipped with an operation
\(\Rightarrow\) satisfying
\[
 c\leq(a\Rightarrow b)
       \quad\Longleftrightarrow\quad c\wedge a\leq b
       \qquad(a,b,c\in H).
\]
This is the residuation law: \(a\Rightarrow b\) is the greatest
element whose meet with \(a\) lies below \(b\).
It determines the operation uniquely.
Write \(\neg a=a\Rightarrow0\).
The one-element algebra, with \(0=1\), is allowed.
A Heyting homomorphism preserves the bounds and all three
binary operations \(\wedge,\vee,\Rightarrow\).
The arrow \(\Rightarrow\) here acts on algebra elements; the
object-language connective is \(\to\), as in
Definition~\ref{def:p08-heyting-algebra}.
The source notation \(a\to b\) for algebra elements denotes this
same uniquely determined residual; thus \(a\to b=a\Rightarrow b\)
when the operands are elements of a Heyting algebra.
\index{Heyting algebra}\index{residuation!Heyting algebra}
\end{hbdefinition}

\begin{hbtheorem}{thm:p10-heyting-examples}{proved=annex:p10.heyting.1}{Open-set and upset Heyting algebras}
For every topological space \(X\), its open subsets form a
Heyting algebra \(\mathcal O(X)\), with finite meets given by
intersection, finite joins by union, bounds \(\varnothing,X\),
and
\[
 U\Rightarrow V=\operatorname{int}_X((X\setminus U)\cup V).
\]
For every preordered set \((W,\leq)\), the upsets form a
Heyting algebra \(\operatorname{Up}(W)\), with the same
set operations and
\[
 U\Rightarrow V
   =\{w\in W:\forall v\geq w\ (v\in U\Longrightarrow v\in V)\}.
\]
These statements include empty \(X\) or \(W\), when the respective
algebra is trivial.  A nonempty \(W\) gives a nontrivial upset
algebra and supports the intuitionistic Kripke semantics of
Part~VIII.
\end{hbtheorem}

The residuation definition appears in
\cite[Definition~2.2.4, p.~21]{Bezhanishvili2006Thesis};
the upset construction and prime-filter order are given there
in \S2.2.3, pp.~26--27.  The proofs below also cover preorders.

\begin{hbdefinition}{def:p10-ipc-lindenbaum}{IPC derivability and its Lindenbaum algebra}
Let \(\mathsf{Prop}\) be any set of propositional variables.
Use the language \(\bot,\land,\lor,\to\) of
Section~\ref{sec:intuitionistic-logic-01}, with
\(\neg\varphi=(\varphi\to\bot)\) and
\(\top=(\bot\to\bot)\).
Write \(\Gamma\vdash_{\mathsf{IPC}}\varphi\) for that
finite-support \(\mathsf{IPL}\) derivability relation.
Thus IPC and IPL denote the same logic, and every proof here
uses the propositional fragment of
\hyperref[sheet:lj]{\texttt{sheet:lj}}.
Uniform substitution transforms a derivation of
\(\Gamma\vdash\varphi\) into one of
\(\sigma[\Gamma]\vdash\sigma(\varphi)\); it does not permit
free substitution in a conclusion while leaving its premises fixed.

Define
\[
 \varphi\sim\psi
 \quad\Longleftrightarrow\quad
 \{\varphi\}\vdash_{\mathsf{IPC}}\psi
       \text{ and }\{\psi\}\vdash_{\mathsf{IPC}}\varphi.
\]
Let \(H_{\mathsf{IPC}}\) have formula classes \([\varphi]\),
with bounds \([\bot],[\top]\) and operations induced by
\(\land,\lor,\to\).
Its order is
\[
 [\varphi]\leq[\psi]
          \quad\Longleftrightarrow\quad
          \{\varphi\}\vdash_{\mathsf{IPC}}\psi.
\]
The next theorem proves that these definitions are well-defined
and give a Heyting algebra.
A valuation \(v:\mathsf{Prop}\to H\) into a Heyting algebra extends
recursively, interpreting the formula connectives by
\(0,\wedge,\vee,\Rightarrow\).
\end{hbdefinition}

\begin{hbtheorem}{thm:p10-heyting-completeness}{proved=annex:p10.heyting.2}{Algebraic completeness of IPC}
The algebra \(H_{\mathsf{IPC}}\) is a nontrivial Heyting algebra.
For every formula \(\varphi\),
\[
 \vdash_{\mathsf{IPC}}\varphi
 \quad\Longleftrightarrow\quad
 [\varphi]=1
 \quad\Longleftrightarrow\quad
 \llbracket\varphi\rrbracket_v^H=1_H
       \text{ for every Heyting algebra }H\text{ and valuation }v.
\]
The canonical valuation \(v_0(p)=[p]\) gives
\(\llbracket\varphi\rrbracket_{v_0}=[\varphi]\), and therefore
refutes every non-theorem.
More generally, for a finite premise set \(\Gamma\),
\[
 \Gamma\vdash_{\mathsf{IPC}}\varphi
 \quad\Longleftrightarrow\quad
 \bigwedge_{\gamma\in\Gamma}\llbracket\gamma\rrbracket_v^H
       \leq\llbracket\varphi\rrbracket_v^H
       \text{ for every }H,v.
\]
The empty meet is \(1\).
\end{hbtheorem}

The theoremhood statement is
\cite[Theorem~2.2.16, p.~25]{Bezhanishvili2006Thesis}.
The annex proves soundness directly for the displayed LJ rules,
then constructs the quotient order and its residual.

\begin{hbdefinition}{def:p10-heyting-prime-frame}{The prime-filter order of a Heyting algebra}
A subset \(P\subseteq H\) is a prime filter when it is a lattice filter containing \(1\),
closed upward and under finite meets, excluding \(0\), and
satisfying
\[
 a\vee b\in P\Longrightarrow a\in P\text{ or }b\in P.
\]
Let \(\operatorname{Pf}(H)\) be all such filters, ordered by inclusion.
For \(a\in H\), put
\[
 \beta_H(a)=\{P\in\operatorname{Pf}(H):a\in P\}.
\]
This is an upset of the prime-filter order.
The filters need not be maximal.
\end{hbdefinition}

\begin{hbtheorem}{thm:p10-heyting-prime-representation}{proved=annex:p10.heyting.3}{Prime-filter representation and the implication witness}
Assume BPIT.  For every Heyting algebra \(H\), the map
\[
 \beta_H:H\longrightarrow
       \operatorname{Up}(\operatorname{Pf}(H),\subseteq)
\]
is an injective Heyting homomorphism.
In particular, for \(a,b\in H\) and \(P\in\operatorname{Pf}(H)\),
\[
 a\Rightarrow b\in P
 \quad\Longleftrightarrow\quad
 \forall Q\supseteq P\
       (Q\in\operatorname{Pf}(H)\text{ and }a\in Q
                     \Longrightarrow b\in Q).
\]
If \(a\Rightarrow b\notin P\), a witnessing prime filter
\(Q\supseteq P\) contains \(a\) and omits \(b\).
If \(H\) is nontrivial, its prime-filter order is nonempty.
For trivial \(H\), it is empty and the displayed map is the
isomorphism between two one-element algebras.
For countable \(H\), these conclusions hold in ZF using the
countable case of Theorem~\ref{thm:p10-prime-filter-separation}.
\end{hbtheorem}

The prime-filter construction is the underlying ordered
representation in
\cite[Theorem~2.3.3 and its proof sketch, p.~28]{Bezhanishvili2006Thesis}.
The image is a Heyting algebra of upsets; it is not asserted
to contain every upset.

\begin{hbtheorem}{thm:p10-ipc-prime-kripke-completeness}{proved=annex:p10.heyting.4}{Strong local Kripke completeness of IPC}
Assume BPIT for an arbitrary variable set.
For every formula set \(\Gamma\) and formula \(\varphi\),
\[
 \Gamma\vdash_{\mathsf{IPC}}\varphi
       \quad\Longleftrightarrow\quad
 \Gamma\entails_{\mathsf{IPL}}\varphi
\]
with the persistent valuations, nonempty preorders, and local
forcing convention of
Definitions~\ref{def:p08-ipl-kripke-structure}--\ref{def:p08-ipl-consequence}.
The same equivalence holds if frames are required to be partial orders.
Every failure of derivability is witnessed on the prime filters
of \(H_{\mathsf{IPC}}\), ordered by inclusion, with
\[
 V(p)=\beta_{H_{\mathsf{IPC}}}([p]).
\]
If the variable set is countable, the entire argument works in
ZF by the countable prime-filter-separation theorem.
\end{hbtheorem}
